\mxnewpage\section*{EV4-T2 -- Préparation EAM}\label{EV4-T2}
\addcontentsline{toc}{section}{EV4-T2 -- Préparation EAM}
\textbf{Sujet d’entraînement}. Ensemble du programme.\par Durée : 2 h. Sans calculatrice. Note sur 20. Entraînements de mai et juin ; EAM de zone à confirmer.
\subsection*{Partie 1 -- Automatismes (6 points)}
Douze questions, \num{0,5} point chacune. Une seule réponse correcte ; aucune justification demandée. Une réponse incorrecte ou absente rapporte 0, sans retrait de points.
\begin{enumerate}
\item Calculer $6+2\times3$.\par A. $24$\quad B. $15$\quad C. $11$\quad D. $12$
\item Calculer $\frac12+\frac13$.\par A. $\frac{5}{6}$\quad B. $\frac{2}{5}$\quad C. $\frac{5}{3}$\quad D. $\frac{1}{6}$
\item Calculer $\dfrac{2^6}{2^{4}}$.\par A. $2$\quad B. $4$\quad C. $1$\quad D. $11$
\item Une valeur de 50 augmente de \pc{20}. Valeur finale ?\par A. $60$\quad B. $10$\quad C. $70$\quad D. $40$
\item Une valeur de 80 est obtenue après une baisse de \pc{20}. Valeur initiale ?\par A. $64$\quad B. $100$\quad C. $96$\quad D. $60$
\item Hausse de \pc{25} puis baisse de \pc{20} : coefficient global ?\par A. $\frac{4}{5}$\quad B. $\frac{21}{20}$\quad C. $1$\quad D. $\frac{9}{10}$
\item La solution de $2x+6=0$ vaut\par A. $-3$\quad B. $12$\quad C. $3$\quad D. $-12$
\item La pente entre $\pt{0;6}$ et $\pt{2;12}$ vaut\par A. $2$\quad B. $-3$\quad C. $3$\quad D. $6$
\item La moyenne de $6, 8, 10$ vaut\par A. $6$\quad B. $24$\quad C. $8$\quad D. $9$
\item Un secteur de $90^\circ$ représente quelle fraction du total ?\par A. $\frac{1}{3}$\quad B. $\frac{1}{2}$\quad C. $\frac{3}{4}$\quad D. $\frac{1}{4}$
\item Une probabilité vaut $\dfrac{2}{5}$. Son complément vaut\par A. $\frac{5}{2}$\quad B. $\frac{2}{5}$\quad C. $\frac{1}{5}$\quad D. $\frac{3}{5}$
\item Pour $f(x)=x^2+6$, $f(-2)$ vaut\par A. $10$\quad B. $8$\quad C. $24$\quad D. $2$
\end{enumerate}
\mxnewpage\subsection*{Partie 2 -- Exercices indépendants (14 points)}
Toute démarche doit être justifiée. Les résultats exacts sont attendus sauf indication contraire.
\mxneed{12}
\subsection*{Exercice 1 -- Un arbre renversé (5 points)}
Deux événements $A$ et $B$ vérifient $P(A)=\dfrac35$, $P_A(B)=\dfrac13$ et $P_{\overline A}(B)=\dfrac34$. On cherche à construire un arbre qui commence par distinguer $B$ et $\overline B$, au lieu de $A$ et $\overline A$.
\begin{enumerate}
\item Calculer $P(A\cap B)$ et $P(B)$. \hfill(1 pt)
\item Calculer $P_B(A)$ et $P_{\overline B}(A)$. \hfill(2 pts)
\item Construire et compléter l’arbre pondéré dont le premier niveau distingue $B$ et $\overline B$, puis le second $A$ et $\overline A$. \hfill(2 pts)
\end{enumerate}
\mxneed{12}
\subsection*{Exercice 2 -- Deux méthodes d’optimisation (5 points)}
La fonction $f$ est définie sur l’intervalle $\intff{1}{6}$ par $f(x)=x(8-x)$. On cherche ses extrema sur ce domaine précis, puis les valeurs de $x$ pour lesquelles $f(x)$ atteint au moins 15.
\begin{enumerate}
\item Donner une forme canonique et un maximum. \hfill(1 pt)
\item Retrouver le maximum par dérivation. \hfill(2 pts)
\item Déterminer le minimum sur le domaine et résoudre $f(x)\geq15$. \hfill(2 pts)
\end{enumerate}
\mxneed{12}
\subsection*{Exercice 3 -- Arguments et symétrie (4 points)}
Dans un repère orthonormé d’origine $O$, les points $A$ et $B$ du cercle trigonométrique sont les images de $\dfrac{7\pi}6$ et de $\dfrac{11\pi}6$. L’angle non orienté $\widehat{AOB}$ sera donné en radians, entre 0 et $\pi$.
\begin{enumerate}
\item Donner leurs coordonnées exactes. \hfill(1 pt)
\item Calculer $\vect{OA}\cdot\vect{OB}$. \hfill(1 pt)
\item Déterminer l’angle non orienté AOB puis AB. \hfill(2 pts)
\end{enumerate}
